\section{Research Questions}
\label{sec:rqs}

\paragraph{RQ1: Uniform utility versus optimized utility.}
Does replacing one attribution update at
$\vomega^{\mathrm{unif}}$ with a persistent sequence of continuous,
budget-feasible weights produce a better Top-$10\%$ subset? We answer with the
paired gap between persistent and reset branches at each common freeze point.
An improvement over one-shot alone is insufficient: persistent optimization
must also beat repeated reset scoring to establish a benefit from outer
memory. The supporting diagnostics are development loss, selected-noise rate,
weight movement, subset turnover, and final performance at the common horizon.

\paragraph{RQ2: How does selection scale with complete bilevel rounds?}
When the full inner-update, outer-score, weight-projection, and Top-$B$
deployment cycle is repeated $R\in\{1,2,3,4,5\}$ times, does performance
improve monotonically or saturate? Every branch trains the target to the same
$H=5L$ horizon, so $R$ changes selection opportunities but not real target
epochs or examples. We report the entire freeze-$R$ curve rather than choosing
one round count on test data, and pair accuracy-versus-target-data with
accuracy-versus-total-compute.

\paragraph{RQ3: Which technical choices create practical RHO's gains?}
We factor RHO into four protocol choices---global versus batch-local quota,
blockwise versus immediate target updates, stale versus online scores, and an
updating joint comparator versus a frozen holdout IL model. At matched protocol
points we also substitute the strict VF, updating-comparator, and frozen-IL
score formulas. This design asks separately whether a gain comes from a more
useful loss difference, greater batch diversity, faster adaptation, or finer
optimizer-step granularity. Batch-size scaling is the priority experiment
because it simultaneously controls the size of the local ranking pool and the
number of immediate target updates.

The algebraic correspondence in Section~\ref{sec:connection} is treated as a
correctness precondition, not a standalone performance RQ. Gradient sign,
inner tracking, budget projection, and ranking agreement must pass the gates in
Section~\ref{sec:experiments} before the three questions above receive a
value-function interpretation.
